% =====================================================================
%  Constantes oficiales de la materia ("Tabla de constantes.pdf").
%  Usar SIEMPRE estos valores en los ejemplos numéricos, para que
%  coincidan con los que se usan en clase y en el examen.
% =====================================================================

% ---------- Constantes físicas ----------
\newcommand{\qval}{\num{1.60e-19}}          % C
\newcommand{\kval}{\num{8.61e-5}}           % eV/K
\newcommand{\kJval}{\num{1.38e-23}}         % J/K
\newcommand{\hval}{\num{6.62e-34}}          % J s
\newcommand{\heVval}{\num{4.13e-15}}        % eV s
\newcommand{\epsval}{\num{8.85e-14}}        % F/cm
\newcommand{\VTval}{\SI{26}{mV}}            % kT/q a 300 K (0.0259 V)
\newcommand{\epsrox}{3.9}                   % SiO2

% ---------- Silicio ----------
\newcommand{\epsrSi}{11.8}
\newcommand{\niSi}{\num{1.0e10}}            % cm^-3
\newcommand{\munSi}{1350}                   % cm^2/(V s)
\newcommand{\mupSi}{480}                    % cm^2/(V s)
\newcommand{\EGSi}{1.12}                    % eV
\newcommand{\NCSi}{\num{3.22e19}}           % cm^-3
\newcommand{\NVSi}{\num{1.83e19}}           % cm^-3

% ---------- Germanio ----------
\newcommand{\epsrGe}{15.8}
\newcommand{\niGe}{\num{2.4e13}}
\newcommand{\munGe}{3900}
\newcommand{\mupGe}{1900}
\newcommand{\EGGe}{0.66}

% ---------- GaAs ----------
\newcommand{\epsrGaAs}{13.1}
\newcommand{\niGaAs}{\num{2.1e6}}
\newcommand{\munGaAs}{8800}
\newcommand{\mupGaAs}{400}
\newcommand{\EGGaAs}{1.42}

% ---------- Tabla para incluir en los apuntes ----------
\newcommand{\tablaconstantes}{%
\begin{tabular}{lccc}
\toprule
 & Si & Ge & GaAs \\
\midrule
$\epsilon_r$                     & 11.8 & 15.8 & 13.1 \\
$\mu_n$ (\si{cm^2/V.s})          & 1350 & 3900 & 8800 \\
$\mu_p$ (\si{cm^2/V.s})          & 480  & 1900 & 400 \\
$n_i$ (\si{cm^{-3}}, 300\,K)     & \num{1.0e10} & \num{2.4e13} & \num{2.1e6} \\
$E_G$ (\si{\eV})                 & 1.12 & 0.66 & 1.42 \\
\bottomrule
\end{tabular}}
